
Formal specification synthesis remains a bottleneck in deploying formal verification to safety-critical systems. When automated verifiers reject specifications, existing approaches discard failure feedback and regenerate from scratch, ignoring structured semantic signals that could guide systematic refinement. This work demonstrates in proof-of-concept that verifier feedback can be viewed as a semantic gradient for specification synthesis. By decomposing feedback into three informational dimensions---Witness Instantiation (concrete counterexamples), Logical Structure (proof obligation categories), and Dependency Preservation (causal chains)---we encode complementary semantic constraints enabling systematic refinement. In a proof-of-concept with simulated verifier feedback, structured multi-dimensional feedback achieves 70.1\% proof discharge rates within 5.3 iterations on average, outperforming unstructured feedback by 38.2 percentage points (70.1\% vs. 31.9\%, $\beta =12.49, R^{2}=0.89$, p<$10^{-50})$ and compute-matched single-shot sampling by 10.7 percentage points (71.4\% vs. 60.8\%, p<0.0001). An 8-primitive semantic normalization layer achieves 100\% error category coverage across verifiers, enabling cross-tool transfer with 15.1\% performance degradation. Mutation testing shows synthesized specifications achieve 63.3\% mutation kill rate, matching expert-written gold baseline (60\%). These proof-of-concept results provide an information-theoretic framework for understanding verification-in-loop.
